\documentclass[letterpaper,11pt]{article}
\usepackage[margin=1.1in]{geometry}

% For angular brackets
\usepackage{textcomp}

% Skip the indents for every paragraph
\usepackage{parskip}
\usepackage{epsfig,endnotes}
\usepackage[T1]{fontenc}
\usepackage{times}
\usepackage[square,comma,numbers,sort&compress]{natbib}
\usepackage[breaklinks,colorlinks]{hyperref}
\hypersetup{citecolor=blue,linkcolor=blue}
\usepackage{listings}
%\usepackage[outputdir=latex.out]{minted}
\usepackage{endnotes,graphicx,fancyvrb,multirow}
\usepackage{booktabs}
\usepackage{underscore}
\usepackage{amssymb}
\usepackage{tikz}
\usetikzlibrary{shapes,arrows,shadows,matrix,positioning,fit}


\newcommand\todo[1]{\textcolor{red}{TODO: #1}}
\newcommand\note[1]{\textcolor{red}{NOTE: #1}}
\newcommand\pf[1]{\textcolor{red}{PF: #1}}

\newcommand\ic[1]{\texttt{#1}}
\newcommand\red[1]{\textcolor{red}{#1}}

\newcommand\para[1]{\vspace*{5pt} \noindent \textbf{#1}}
\newcommand\TODO{\textcolor{red}{TODO}}

\newenvironment{code}
{\begin{center}
	\fontfamily{pcr} \selectfont
    \begin{tabular}{|p{0.9\textwidth}|}
    \hline
	\\[-8pt]
    }
    {
	\\[-8pt]
	\\
	\hline
    \end{tabular}
    \end{center}
}

\def\Snospace~{\S{}}
\renewcommand*\sectionautorefname{\Snospace}
\def\sectionautorefname{\Snospace}
\def\subsectionautorefname{\Snospace}
\def\subsubsectionautorefname{\Snospace}


\newenvironment{temporary}{\par\color{blue}}{\par}


\newcommand{\squishenumerate}{
   \begin{enumerate}
    { \setlength{\itemsep}{0pt}      \setlength{\parsep}{3pt}
      \setlength{\topsep}{3pt}       \setlength{\partopsep}{0pt}
      \setlength{\leftmargin}{2em} \setlength{\labelwidth}{1em}
      \setlength{\labelsep}{0.5em} } }

\newcommand{\squishenumerateend}{
    \end{enumerate}  }

\newcommand{\squishlist}{
   \begin{list}{$\bullet$}
    { \setlength{\itemsep}{0pt}      \setlength{\parsep}{3pt}
      \setlength{\topsep}{3pt}       \setlength{\partopsep}{0pt}
      \setlength{\leftmargin}{2em} \setlength{\labelwidth}{1em}
      \setlength{\labelsep}{0.5em} } }

\newcommand{\squishlisttwo}{
   \begin{list}{$\bullet$}
    { \setlength{\itemsep}{0pt}    \setlength{\parsep}{0pt}
      \setlength{\topsep}{0pt}     \setlength{\partopsep}{0pt}
      \setlength{\leftmargin}{2em} \setlength{\labelwidth}{1.5em}
      \setlength{\labelsep}{0.5em} } }

\newcommand{\squishend}{
    \end{list}  }


% For customizing itemize spacing
\usepackage{enumitem}
\setlist[itemize]{topsep=0pt, itemsep=0pt, parsep=0pt}
\setlist[enumerate]{topsep=0pt, itemsep=0pt, parsep=0pt}

% For those sweet, sweet acmart fonts
\usepackage{inconsolata}
\RequirePackage[T1]{fontenc} \RequirePackage[tt=false,type1=true]{libertine}

% For more colors!
\usepackage[dvipsnames]{xcolor}

\begin{document}

\section*{\textcolor{BlueViolet}{Course Information}}

\begin{itemize}
\item \Large{\textbf{Name:} CS 408: Software Testing}
\item \textbf{Course Number:} CS 40800 - LE1
\item \textbf{Course Reference Number/CRN:} 13463
\end{itemize}

\section*{\textcolor{BlueViolet}{General Course Description}}
This course covers the following topics:
\begin{itemize}
\item errors and testing;
\item software quality, requirements, behavior, and correctness;
\item testing, debugging, verification;
\item control flow graphs, dominators;
\item types of testing;
\item Test selection: from requirements, finite state models, and combinatorial
  designs;
\item regression testing and test minimization;
\item Test adequacy
  assessment: control and data flow;
\item mutation based;
\item testing tools
\end{itemize}

\section*{\textcolor{BlueViolet}{Course Learning Outcomes}}
After successfully completing this course, students will:

\begin{itemize}

\item \textbf{Become better testers.}
Understand and be able to create high-quality tests at all testing levels.
Learn practical ways to design and automate high-quality tests.
Understand how to apply theory (test criteria) in practical ways.

\item \textbf{Become better programmers.}
Be aware of potential problems in software and able to create high-quality developer tests.

\item \textbf{Become better engineers.}
Understand and be able to build programs and test them in a unified manner.

\item \textbf{Become better thinkers.}
Understand and be able to approach software problem-solving in logical and analytical ways.
\end{itemize}

\section*{\textcolor{BlueViolet}{Learning Resources / Text}}
The main textbook for this course will be:
\begin{itemize}
\item \textit{Introduction to Software Testing}, 2nd edition, by Paul
  Ammann and Jeff Offutt.
\end{itemize}
Supplementary material will be drawn from a variety of sources, and
will be posted to the course discussion board.

\section*{\textcolor{BlueViolet}{Types of Assignments/Assessments}}
The course features three categories of assessments:
\begin{itemize}
  \item \textbf{Homeworks:} These are pen and paper assignments that
    students need to complete individually.
  \item \textbf{Group Projects::} Programming projects that will be
    done in groups.
  \item \textbf{Exams:} There will be an in-class / in-person midterm
    roughly halfway through the semester, and another in-person final.
  \end{itemize}

\section*{\textcolor{BlueViolet}{Grading Scale}}
The weighted contribution of each category of assessment to final
grades will be:
\begin{itemize}
\item Projects: 25\%
\item Homework: 15\%
\item Midterm (in-person): 30\%
\item Final (in-person): 30\%
\end{itemize}

At the end of the semester, final grades will be calculated by summing
up your weighted homework and test grades and translating the
truncated result into letter grades using the following buckets as an
\emph{upper bound}, i.e., the thresholds for each grade category may
be lower. The thresholds for + and – will be determined at the end of
the semester:
\begin{itemize}
\item \textbf{A}: 100-90
\item \textbf{B}: 89-80
\item \textbf{C}: 79 - 70
\item \textbf{D}: 69 - 60
\item \textbf{F}: 59 and below
\end{itemize}

\end{document}
